\uuid{xKnc}
\exo7id{3501}
\titre{exo7 3501}
\auteur{quercia}
\organisation{exo7}
\datecreate{2010-03-10}
\isIndication{false}
\isCorrection{true}
\chapitre{Réduction d'endomorphisme, polynôme annulateur}
\sousChapitre{Valeur propre, vecteur propre}
\cpge{mpsi-mp2i.determinants.e, pcsi.determinants.d, mpsi-mp2i.calcul-algebrique-trigonometrie.d, pcsi.calcul-algebrique-trigonometrie.d, ptsi.calcul-algebrique-trigonometrie.d, bcpst1.trigonometrie, mp-mpi.reduction-endomorphismes.c, pc.reduction.b, psi.reduction-endomorphismes.b, pt.reduction-endomorphismes.b, mp-mpi.reduction-endomorphismes.b, pc.reduction.a, psi.reduction-endomorphismes.a, pt.reduction-endomorphismes.a, bcpst2.valeurs-propres-vecteurs-propres.a}
\module{Algèbre}
\niveau{L2}
\difficulte{}

\contenu{
\texte{
Soit $A = \begin{pmatrix} 0 &1      &       &(0)     \cr
                    1 &\ddots &\ddots          \cr
                      &\ddots &\ddots &1       \cr
                    (0)&      &1      &0       \cr \end{pmatrix} \in \mathcal{M}_n(\R)$.
}
\begin{enumerate}
    \item \question{Calculer $D_n(\theta) = \det(A+(2\cos\theta) I)$ par récurrence.}
\reponse{$D_n = 2\cos\theta D_{n-1} - D_{n-2}  \Rightarrow 
              D_n = \frac{\sin(n+1)\theta}{\sin\theta}$.}
    \item \question{En déduire les valeurs propres de $A$.}
\reponse{$-2\cos\left(\frac{k\pi}{n+1}\right)$, $1 \le k \le n$.}
\end{enumerate}
}
